\begin{center}
    \large
    
    \begin{singlespace}
        \textbf{\englishTitle{}} \\[0.5cm]
    \end{singlespace}
    
    \begin{singlespace}
        Student : \studentEnName{}  \hspace{1.0cm} 
        % 兩個指導教授要寫 Advisors
        \ifdefined\advisorCnNameB
            Advisors: Dr.\, \advisorEnName \\
            \hspace{6.6cm} Dr.\, \advisorEnNameB  \\
        \else
            Advisor: Dr.\, \advisorEnName \\
        \fi
    \end{singlespace}
    
    \vspace{0.5cm}
    \begin{singlespace}
        \NameofDepartmentInstituteEN\\
        National Yang Ming Chiao Tung University\\[0.2cm]
    \end{singlespace}
    
    \textbf{Abstract} \\[0.5cm]

\end{center}

\normalsize

Valley excitons in transition metal dichalcogenides (TMDs) naturally carry spin angular momentum (SAM). Meanwhile, orbital angular momentum (OAM) can theoretically provide unlimited degrees of freedom. Therefore, converting valley SAM into OAM selectively has become an important topic for developing high-capacity quantum light sources and valley photonics applications.

This study proposes an inverse design framework that combines the 3D finite-difference time-domain (FDTD) method with the adjoint method. Using this framework, we designed a high-degree-of-freedom gallium phosphide ($\text{GaP}$) topology-optimized structure. Within a very small fixed volume, this structure can deterministically convert the valley SAM of monolayer tungsten disulfide ($\text{WS}_2$) excitons into OAM with any designated topological charge, producing high-purity Laguerre-Gaussian (LG) beams.

Numerical simulations for a full series from $\ell = 0$ to $5$ prove this design. At the emission wavelength of $\text{WS}_2$ excitons, the structure was optimized for left-circularly polarized (LCP) emission. The far-field mode purity of the final device reached over $80\%$. In addition, when the exciton emission switched to right-circularly polarized (RCP), the far-field purity of the target LG mode drops below $1\%$, showing very high valley polarization selectivity. These results provide a new strategy for designing miniaturized vortex optical and valleytronic devices.
\vspace{1cm}

\textbf{Keywords: } Valleytronics, Transition Metal Dichalcogenides, Laguerre-Gaussian beam, Metasurface, Topology Optimization, Adjoint Method.